% ============================================================
% SOLUTION — PART III: FILL IN THE BLANKS
%
% AI INSTRUCTION:
% - Reprint each sentence with the blanks FILLED IN and highlighted
%   using \correct{...}, so the student reads a complete correct sentence
% - Where several wordings are acceptable, list them:
%   "accepted: X, Y, or Z"
% ============================================================

\section*{Part III: Fill in the Blanks --- Solution}

\textbf{3.1}
\begin{ans}
A sentence with one or more blanks embedded in a definition:
the \correct{answer} property guarantees that \correct{answer}
even when the system is under load.

\textit{Also accepted:} alternative wording that shows the same understanding.
\end{ans}

\textbf{3.2}
\begin{ans}
% \correct{...} is a TEXT-mode macro: wrap math inside it, never the reverse.
$b = $ \correct{$\mathbf{As} + \mathbf{e}$} $\pmod{q}$; the added term
\correct{$\mathbf{e}$ (the error/noise vector)} is what makes the system hard.
\end{ans}

\begin{keypoint}
The single fact this whole question was testing.
\end{keypoint}

\textbf{3.3}
\begin{ans}
\correct{quadratic} speedup, reducing the cost to approximately
\correct{$2^{64}$} operations.
\end{ans}

\textbf{3.4}
\begin{ans}
Answer derived from the table: \correct{the filled-in answer}.

Show the one-line reasoning that reads the answer off the table.
\end{ans}

\textbf{3.5}
\begin{ans}
XXXX stands for \correct{the full expansion} and is used to
\correct{its purpose}.
\end{ans}

\textbf{3.6}
\begin{ans}
In order: \correct{step one}, \correct{step two}, \correct{step three}.
\end{ans}

\begin{caution}
Order matters here --- listing the right three items in the wrong order
earns partial credit only.
\end{caution}
